\documentclass[11pt,twoside]{amsart}
\usepackage{amssymb, amsmath, enumerate, palatino, hyperref}
\usepackage{caption}
\usepackage[normalem]{ulem}
\usepackage{fullpage}
\usepackage{tikz}
\usetikzlibrary{arrows.meta,positioning}

\usepackage[T1]{fontenc}
\renewcommand{\labelitemi}{\guillemotright}
\usepackage{mathrsfs}

\theoremstyle{plain}
\newtheorem{prop}{Proposition}
\newtheorem{lemma}[prop]{Lemma}
\newtheorem{thm}[prop]{Theorem}
\newtheorem{obs}[prop]{Observation}
\newtheorem{app}[prop]{Application}
\newtheorem*{MainThm}{Main Theorem}
\newtheorem{cor}[prop]{Corollary}
\newtheorem{conj}[prop]{Conjecture}
\theoremstyle{remark}
\newtheorem{rmk}[prop]{Remark}
\newtheorem*{rmk*}{Remark}
\newtheorem{prob}{Problem}
\newtheorem*{bonus*}{Bonus Problem}
\newtheorem{exc}{Exercise}
\theoremstyle{definition}
\newtheorem{ex}[prop]{Example}
\theoremstyle{definition}
\newtheorem{defn}[prop]{Definition}

\newcommand{\RR}{\mathbb{R}}
\newcommand{\ZZ}{\mathbb{Z}}
\newcommand{\CC}{\mathbb{C}}
\newcommand{\NN}{\mathbb{N}}
\newcommand{\QQ}{\mathbb{Q}}
\newcommand{\FF}{\mathbb{F}}

\newcommand{\tr}{\operatorname{tr}}
\newcommand{\defeq}{:=}
\newcommand{\sm}{\smallsetminus}
\newcommand{\lra}{\longrightarrow}
\newcommand{\ol}[1]{\overline{#1}}
\newcommand{\spec}{\operatorname{spec}}
\newcommand{\Awalks}[1]{w_{#1}}

\title{Math 372: Combinatorics\\ Homework 01}

\author{Due Friday, 11 September, 10\textsc{p.m.}, via Gradescope}

\begin{document}

\maketitle

\begin{rmk*}
\textbf{Conventions.} All graphs here are finite and undirected, and they are allowed to have
loops. If $G$ has vertex set $\{v_1,\dots,v_p\}$, its \emph{adjacency matrix} $A(G)$ is the
$p\times p$ matrix whose $(i,j)$ entry is the number of edges joining $v_i$ to $v_j$; a loop at
$v_i$ contributes $1$ to the $(i,i)$ entry. A \emph{walk of length $\ell$} from $u$ to $v$ is a
sequence $u=u_0,u_1,\dots,u_\ell=v$ of vertices in which each consecutive pair is joined by an
edge, and the walk is \emph{closed} if $u=v$. Since $A(G)$ is a real symmetric matrix, its
eigenvalues are real; the \emph{spectrum} $\spec(G)$ is the multiset of those eigenvalues, listed
with multiplicity. Throughout, $\ell\ge 1$.
\end{rmk*}

\begin{rmk*}
\textbf{On these solutions.} Each problem is worth up to five points for mathematical content
and up to two points for writing. Write your solutions as explanations addressed to another
student in this course: say what you are going to show before you show it, justify each step,
and use complete sentences even when symbols are doing most of the work. End your submission
with a short \emph{Methods and tools} note --- three or four sentences saying what software you
used and for what, and how you checked anything a machine produced for you.
\end{rmk*}

%%% Problem 1 %%%
\begin{prob}
Let $D$ be the graph on four vertices pictured below, with adjacency matrix $A$ taken with
respect to the ordering $a,b,c,d$:
\begin{center}
\begin{tikzpicture}[every node/.style={circle,draw,inner sep=1.6pt,fill=black}]
  \node[label=left:$a$]  (a) at (0,1) {};
  \node[label=right:$b$] (b) at (2,1) {};
  \node[label=left:$c$]  (c) at (0,0) {};
  \node[label=right:$d$] (d) at (2,0) {};
  \draw (a)--(b); \draw (a)--(c); \draw (a)--(d);
  \draw (b)--(c); \draw (b)--(d);
\end{tikzpicture}
\qquad
$A=\begin{pmatrix} 0&1&1&1\\ 1&0&1&1\\ 1&1&0&0\\ 1&1&0&0 \end{pmatrix}$
\end{center}
\begin{enumerate}[(a)]
\item Determine $\spec(D)$.
\item Determine the number of closed walks of length $4$ in $D$ in two genuinely different ways:
once using your answer to (a), and once by a computation that does not use the eigenvalues at
all. Confirm that the two answers agree.
\end{enumerate}
The mathematics here is routine, and that is deliberate. This problem exists so that we can agree
on what a written solution in this course looks like, so spend your effort on the exposition: a
reader who has not seen $D$ should be able to follow your work from beginning to end without
supplying anything themselves.
\end{prob}

%%% Problem 2 %%%
\begin{prob}
Suppose $G$ is a graph with $15$ vertices, and suppose that for every $\ell\ge 1$ the number of
closed walks of length $\ell$ in $G$ is
\[
8^{\ell} \;+\; 2\cdot 3^{\ell} \;+\; 3\cdot(-1)^{\ell} \;+\; (-6)^{\ell} \;+\; 5^{\ell}.
\]
Let $G'$ be the graph obtained from $G$ by adding a loop at every vertex, in addition to whatever
loops $G$ already has. How many closed walks of length $\ell$ does $G'$ have?

\smallskip
Give a complete justification, not merely the formula. (Linear-algebraic reasoning is the
intended route. If you would also like to think about what the answer says combinatorially, you
are welcome to, but it is not required.)
\end{prob}

%%% Problem 3 %%%
\begin{prob}
Let $r,s\ge 1$. The \emph{complete bipartite graph} $K_{r,s}$ has vertices
$u_1,\dots,u_r,v_1,\dots,v_s$, with exactly one edge joining $u_i$ to $v_j$ for every $i$ and $j$,
and no other edges. (So $K_{r,s}$ has $r+s$ vertices and $rs$ edges.)
\begin{enumerate}[(a)]
\item By purely combinatorial reasoning --- that is, by counting walks directly, without
computing any eigenvalues --- determine the number of closed walks of length $\ell$ in $K_{r,s}$.
\item Deduce $\spec(K_{r,s})$ from your answer to (a).
\end{enumerate}
Part (b) runs in the opposite direction from everything we did in class this week: there we
computed a spectrum and read off walk counts, and here you are asked to compute walk counts and
read off a spectrum. Say clearly what makes that deduction valid.
\end{prob}

%%% Problem 4 %%%
\begin{prob}
A graph $G$ is \emph{bipartite} if its vertex set can be written as a disjoint union
$V = X \sqcup Y$ in such a way that every edge of $G$ joins a vertex of $X$ to a vertex of $Y$.
(In particular a bipartite graph has no loops.)

\smallskip
Show that if $G$ is bipartite, then its nonzero eigenvalues come in pairs $\pm\lambda$: that is,
for each $\lambda\neq 0$, the eigenvalue $\lambda$ and the eigenvalue $-\lambda$ occur in
$\spec(G)$ with the same multiplicity.

\smallskip
Give a \emph{walk-counting} argument. There is a slicker proof available using eigenvectors, and
you are welcome to find it and include it as a remark, but the argument you submit should proceed
by counting walks.
\end{prob}

%%% Bonus %%%
\begin{bonus*}
Let $G$ be a graph and let $\Delta$ be the largest degree of any vertex of $G$, where the
\emph{degree} of $v$ is the number of edges incident to $v$. Let $\lambda_1$ be the largest
eigenvalue of $A(G)$. Show that $\lambda_1 \le \Delta$.
\end{bonus*}

\end{document}
